\documentclass[12pt]{article}
\usepackage[utf8]{inputenc}
\usepackage[T1]{fontenc}
\usepackage{hyperref}
\title{Sylheti Language Resources for K-12 Curriculum Development}
\author{Compiled by Claude Sonnet 5.5 (claude-sonnet-5-5) for Andr\'as Kornai}
\date{}
\begin{document}
\maketitle

\section{Introduction}
Sylheti (ISO 639-3 syl, Glottocode sylh1242) is an Indo-European language from the Sylhet region of Bangladesh and the Barak Valley of India, with a large diaspora. The reference description gives about 100 speakers, which is clearly wrong. The search results give estimates of 10 to 12 million speakers, though sources disagree. It is written in Syloti Nagri, Eastern Nagari and Latin scripts. Many of the resources below are community-run and should be checked before use.

\section{Primary Educational Resources}
\begin{itemize}
\item Let's Learn Sylheti: spoken lessons, writing tutorials, book recommendations and a download library. \url{https://sylhetilanguage.net/learn-sylheti/}
\item Wikibooks Sylheti course pages, which were mostly outlines. \url{https://en.wikibooks.org/wiki/Sylheti/First_Person} and \url{https://en.wikibooks.org/wiki/Sylheti/Third_Person}
\end{itemize}

\section{Linguistic and Reference Materials}
\begin{itemize}
\item Sylheti Language Dictionary, a community project with pronunciation, transliteration, meanings and examples in Syloti Nagri. \url{https://sylhetilanguage.net/test/}
\item Living Dictionaries page for Syloti. \url{https://livingdictionaries.app/syloti/about}
\item Omniglot overview of the language and its script. \url{https://omniglot.com/writing/sylheti.htm}
\item Wiktionary entries for Syloti Nagri characters. \url{https://en.wiktionary.org/wiki/%EA%A0%8B} and \url{https://en.wiktionary.org/wiki/%EA%A0%86}
\item A page of Sylheti language resources. \url{https://universeofmemory.com/sylheti-language-resources/}
\end{itemize}

\section{Religious and Community Resources}
None specific were found. The Linguists Collective is mentioned on the Let's Learn Sylheti site as a source of tuition and support.

\section{Digital and Audio Resources}
Let's Learn Sylheti has spoken lessons. The search also mentioned Memrise, Quizlet and YouTube playlists, but no URLs were seen for them. The Endangered Alphabets page covers the script. \url{https://www.endangeredalphabets.net/?p=408}

\section{Curriculum Development Context}
Sylheti has a large speaker base and a historic script, with active community resources. No school programme was found. Sylheti is often treated as a dialect of Bengali, so schooling is normally in Bangla. A curriculum would build on the community dictionary and lessons. It would need decisions on script (Syloti Nagri versus Latin or Bengali script) and on academic verification of grammar.
\end{document}